\section{九〇七至九六〇：逐年军费、区域与社会}
\label{sec:five-annual-military-fiscal-regional-society}

唐亡至北宋建立的编年，不能只由五次中原改朝换代串联。军队要有粮饷、城市要有市场和仓储、政权要能任命州县与处理契约；这些条件在汴洛、河东、燕云、两浙、福建、岭南与四川的节奏不同。正史和《资治通鉴》可核定许多即位、战役和任免，钱币、墓志、寺院碑刻、地方志与考古则使区域生活留下片段。它们都不构成一份完整的人口或财政总账。

\subsection{九〇七至九一二：唐亡以后没有同一天的断裂}

九〇七年朱温建立后梁，首先改变的是中原的朝廷名义和汴州周围的军政网络。诏令、年号和任命可以确证，县乡的赋役、市场和家族关系却不会在同日转换。后梁依赖运河、粮道、盐利与军队，说明军事权力具有物质基础；不能从汴梁朝廷的文书推定河北、淮南或江南均已按同一方式纳入。

九〇八年李克用去世、李存勖继承河东，九〇九年至九一二年河东与后梁交战，城池、关隘和军粮反复成为年表节点。战报中的兵数、斩获和居民态度服务胜败叙事，应与不同政权的记录互校。河东军人、当地农户、商旅和被迁徙者的处境不相同；“沙陀”在史书中的用法也不能替代每个人的自称或社会经验。

\subsection{九一三至九二三：战争、运河与地方中介}

九一三年后梁内部政变、九一四至九二一年河朔争夺，显示继承危机与军队支付相互牵连。兵变不能只以忠奸解释；赏赐、欠饷、将领亲属、驻地市场和道路安全都构成可能条件。史料较少保存普通军户的家庭账，只能确认军费与军纪的联系，不能为所有士卒指定同一种动机。

九二三年后唐灭后梁，是中原政权转换的重要年点。接收旧都意味着重开仓库、安排旧官、确认税源并处理降军，绝非只改年号。官方纪事强调新王朝的合法性，地方恢复速度却受道路、粮食和地方强人的影响。后唐的存在不应被写成对所有区域的一次性统一。

\subsection{九二四至九三六：宫廷政变与区域财政}

九二四年至九二六年后唐的宫廷矛盾与兵变，显示中枢权力依赖军队、俸给和运输。九二六年李嗣源即位可由年序明确，却不说明每一军镇自动改隶。官员传记偏重决策者，军粮文书和地方材料才可能显示某些地区如何承担征发；保存的不均衡使我们不能计算全国军费。

九二七年至九三二年，吴、吴越、闽、南汉与前后蜀分别经营水路、盐货、茶货、城市与寺院。南方政权并不是中原战乱的静止背景。钱币、港口遗址、碑刻和文书可以证明局部交换、供养或官制，却不能将港市繁荣换算为全境富裕，亦不能把王室赞助等同于居民政治忠诚。

九三三年至九三六年石敬瑭与河东、契丹关系的发展，预示新的中原军事联盟。外交文本、使节活动和军事行动是不同层次：盟约能够说明承诺的表达，不能自动证明关隘居民、牧民、商旅或军队将按文书安排生活。边界应被理解为道路、市场和驻军不断协商的地带，而非静止线条。

\subsection{九三六至九四七：后晋、燕云与北方社会}

九三六年后晋建立并向契丹求援，燕云十六州问题由此进入中原政权的长期约束。割让、驻军、税粮和迁徙必须分别讨论：一份政治承诺不等于所有地方行政立即重组。汉文正史、契丹材料和边地遗存对事件的命名不同，阅读时应保留其政治立场，避免把其中任一方的“归属”叙述当作全体居民的经验。

九四〇年前后后晋的财政、盐铁和军政措施，面对的是各镇、军户和市场的差异。制度条文说明朝廷希望征收什么，不能证明应收已成实收，亦不能由中央岁入推断家庭负担均等。货币、绢帛、粮食与劳役在不同地区可互相折算或并存；“货币化”不能替代对具体税种和支付方式的考察。

九四七年契丹入汴、后晋覆亡，随后后汉建立。首都易手造成仓储、官员、城市治安和交通的危机，但并非所有省份同时陷入同一秩序。流民和降军的数字多出自军政叙事，适合显示动员规模的表达而非精确人口统计。地方社会的恢复仍须追踪户籍、土地、寺院与市场在何处重新连接。

\subsection{九四八至九五一年：后汉与后周的接收}

九四八至九五〇年后汉内部权力变化表明短命王朝仍需面对军人拥立、饷源和州镇关系。所谓“军阀”不是一种脱离行政的个人力量；他们必须通过税粮、商路、地方精英和文书网络维持部众。人物传记可观察联盟，不能替代州县社会对征发与保护的多样反应。

九五一年郭威建后周，新的政权接收后汉官员、仓储、军队与旧有债务关系。接收政策可以从诏令和任命中确认，实际执行会因地方官、更替速度和资源而异。把后周改革描述为中央能力突然恢复，容易遮蔽战争年代一直存在的地方中介和区域差别。

\subsection{九五二至九五六：货币、寺院与军队}

九五二至九五四年后周整顿军队、州县与财政，军籍、军粮和家庭劳作相互牵连。军队不是只在战场出现的集体：驻地消费、运输、婚姻和土地使用都可能改变城市与乡村。现存资料不能充分记录普通军人的家属，叙事不应假定他们都拥有同样身份或收益。

九五四年高平之战后，军政资源向北方防务集中。战役结果可以由多种编年材料钉住，伤亡和战功数字则需保留战报修辞。战争既影响中央合法性，也会重排道路、马匹、粮仓和劳力；但一场胜利不足以说明边地行政已稳固。

九五五年至九五六年围绕寺院、僧尼、铜像与铸钱的政策，体现宗教、财政和物质资源的交会。诏令和后出的制度记载可以说明政府意图，不能证明所有地方寺院受到同样清查，亦不能以还俗人数断言信仰消失。题记、遗址和僧传所见的恢复与供养也只是局部，不应反向证明政策毫无效果。

\subsection{九五七至九六〇：征淮南与新朝的门槛}

九五七年至九五九年后周南征淮南，使江淮水路、城镇、仓储和地方军队成为军事财政问题。占城、受降和设官各有不同的证明范围：战报能说明军队宣称控制的地点，不等于乡村税役和日常司法已经稳定。淮南居民、盐商、船户和旧军人的处境不能由一个“归附”词概括。

九五九年世宗去世、幼主继位后，中枢军政的脆弱性重新显现。九六〇年赵匡胤即位是后周年序的终点，却不意味着吴越、南唐、荆南、北汉等区域的市场、官制和宗教网络立即改变。北宋前夜应写作不同地区同时面对的政治消息、道路变化与征发预期，而非一幅已经完成的统一地图。

\subsection{年度叙事的材料限度}

逐年列出政变、战役、税制和宗教政策，容易给人一种国家能看清每一处社会变化的错觉。正史和编年史保留的多是中枢与军事节点；寺院碑刻、墓志、钱币、文书和遗址各自只打开局部窗口。它们的价值在于使军费、区域社会和信仰实践不被王朝年号吞没，而不是让写作者为沉默的大多数编造精确的收入、身份或动机。
